\SourceLeaf{347}%
\par\indent Bei Untersuchungen, in denen viele Richtungen gerader Linien im Raum eine Rolle spielen, bezeichnet man diese Richtungen zweckmäßig durch Punkte auf der Oberfläche einer festen Kugel. Es sind die Endpunkte der Radien, die zu den jeweiligen Geraden parallel gezogen werden. Mittelpunkt und Radius dieser \textit{Hilfskugel} sind ganz beliebig. Als Radius kann man die Längeneinheit wählen. Dieses Verfahren stimmt im Wesentlichen mit dem Verfahren überein, das in der Astronomie ständig angewandt wird. Dort bezieht man alle Richtungen auf eine gedachte Himmelskugel mit unendlich großem Radius. Die sphärische Trigonometrie und einige weitere Lehrsätze dienen dann zur Lösung der Aufgaben, die sich beim Vergleich der verschiedenen Richtungen ergeben können. Der Verfasser hat diesen Lehrsätzen noch einen neuen hinzugefügt, der sich häufig anwenden lässt.

Bezeichnet man die Richtung der in jedem Punkt einer gekrümmten Fläche auf ihr errichteten Normalen nach dem beschriebenen Verfahren durch den entsprechenden Punkt auf der Kugeloberfläche, so ordnet man jedem Punkt der gekrümmten Fläche einen Punkt auf der Oberfläche der Hilfskugel zu. Dann entspricht im Allgemeinen jeder Linie auf der gekrümmten Fläche eine Linie auf der Hilfskugel und jedem Flächenstück der einen ein Flächenstück der anderen. Je weniger das Flächenstück von einer Ebene abweicht, desto kleiner ist der entsprechende Teil der Kugeloberfläche. Es liegt daher nahe, den Flächeninhalt des entsprechenden Kugelstücks als Maß für die Totalkrümmung des gekrümmten Flächenstücks zu verwenden. Der Verfasser nennt diesen Flächeninhalt daher die \textit{ganze Krümmung} des entsprechenden Stücks der gekrümmten Fläche. Neben der Größe kommt auch die \textit{Lage} der Teile in Betracht. Ganz unabhängig von ihrem Größenverhältnis kann die Lage in den beiden Flächenstücken ähnlich oder umgekehrt sein. Ein positives oder negatives Vorzeichen der Totalkrümmung kann diese beiden Fälle unterscheiden. Diese Unterscheidung hat jedoch nur dann eine bestimmte Bedeutung, wenn man die Figuren auf bestimmten Seiten der beiden Flächen betrachtet. Der Verfasser nimmt bei der Kugelfläche die Außenseite und bei der gekrümmten Fläche diejenige Seite, auf der man sich die Normale errichtet denkt. Daraus folgt, dass bei konvex-konvexen oder konkav-konkaven Flächen, die sich nicht wesentlich unterscheiden, das positive Vorzeichen gilt und %
